% META: {"id":"1SPE-DERGLOBAL-EX-052","chapitre":"1SPE-DERIVATION-GLOBAL","type_objet":"exercice","capacites":["1SPE-DERIVATION-GLOBAL-C1"],"capacites_codes":["C1"],"methodes":["M1"],"parcours":2,"competences":["raisonner","representer"],"duree_min":10,"mode_creation":"ex_nihilo","sources_inspiration":[],"fichier_tex":"chapitres/1SPE-DERIVATION-GLOBAL/exercices/1SPE-DERGLOBAL-EX-052.tex","corrige_tex":"chapitres/1SPE-DERIVATION-GLOBAL/corriges/1SPE-DERGLOBAL-CO-052.tex","status":"approved","origin":"generated","status_history":["generated","audited","accepted","approved"],"release_acceptance":"RELEASE_OWNER_BATCH_ACCEPTANCE","acceptance_closure_digest":"sha256:9b3ccf9a81c5520fbb7e03b7b2d3e7bf2057a02834b6d908bf3be5f49f3b553f","acceptance_receipt_digest":"sha256:4fad86f0282e0fa4e0419cca1ad6379ae7af5a280c04a4a90880f90a1c044242"}
% BEGIN-VERIFY
% from sympy import symbols, Abs, Rational, limit
% x, h = symbols('x h')
% # Taux de variation de |x| en 0
% taux = (Abs(0 + h) - Abs(0)) / h
% # Pour h > 0 : |h|/h = 1
% assert Abs(1)/1 == 1
% # Pour h < 0 : |h|/h = -1
% assert Abs(-1)/(-1) == -1
% # f(x)=|x-2| : taux de variation en x=2
% # Pour h > 0 : |h|/h = 1, pour h < 0 : |h|/h = -1
% assert Abs(2) == 2
% END-VERIFY
\begin{exercice}{1SPE-DERGLOBAL-EX-052}{2}{10}
On considère la fonction $f$ définie sur $\mathbb{R}$ par $f(x) = |x|$.
\begin{enumerate}
  \item Calculer le taux de variation de $f$ entre $0$ et $h$, pour $h
\neq 0$. Distinguer les cas $h > 0$ et $h < 0$.
  \item En déduire que $f$ n'est pas dérivable en $0$.
  \item La fonction $g$ définie par $g(x) = |x - 2|$ est-elle dérivable en $x = 2$ ? Justifier.
  \item Donner l'expression de $f'(x)$ pour $x > 0$ et pour $x < 0$.
\end{enumerate}
\end{exercice}
